\chapter{Fazit}
Insgesamt lieferten beide Teile des Verusuchs Ergebnisse, die mit der Hypothese übereinstimmten. Zunächst wurde die spektrale Verteilung der kosmischen Strahlung bestimmt, wobei eine Korrelation in $cos^2(\theta)$ festgestellt wurde. Der Exponent betrug:
\begin{align*}
    c = (2{,}01 \pm 0{,}09)
\end{align*}
Der ermittelte Exponent stimmt innerhalb $1\sigma$-Umgebung mit dem theoretisch erwarteten Wert von $c = 2$ ein.
Dieses Ergebnis bestätigt, das die Winkelverteilung die Form $\cos^2\theta$ hat.
Die Anzahl der zufälligen Koinzidenzen war ebenfalls konsistent mit den theoretischen Erwartungen: es wurden mehr räumliche Koinzidenzen (N=1210) als zufällige zeitliche (N=0) Koinzidenzen registriert, was auf eine korrekte Nutzung der Koinzidenzlogik hindeutet.

Das aufgezeichnete Pulshöhenspektrum entspricht eine Landau-Verteilung, ähnlich der von dünnen Szintillatoren. Die Landau-Anpassung ergab:
\begin{align*}
    \chi^2_{\text{red}}=1,35
\end{align*}
besser als bei der Gauß-Verteilung mit dem Wert $\chi^2_{\text{red}}=3,13$. Dies bestätigt, dass die Landauverteilung ein passendes Model für die experimentellen Werte ist und entspricht der Erwartung für dünne Szintillatoren, bei denen das Pulshöhenspektrum durch eine Landau-Verteilung beschrieben wird.

Im zweiten Teil des Versuchs wird die Lebensdauer der Myonen bestimmt. Der erhaltene Wert
$\tau = (\num{0.631} \pm \num{0.001})\,\si{\micro\second}$ \\
Dieser Wert liegt deutlich unter dem Literaturwert von \(\tau = \SI{2.1969811(22)}{\micro\second}\) und passt mit nicht mit diesem innerhalb einer $3\sigma$-Umgebung überrein.
Die während des Veruchs identifizierten systematischen Unsicherheiten und Fehlerquellen, insbesondere jene, die mit der Triggerlogik, Hintergrundereignissen und statistischen Daten zusammenhängen; erklären diesen Unterschied. 
